% ID: SEL-D0045-OB-E01
% ESTADO: validada. S046, cotejada con examen y criterios.
\textbf{a)} $2(C-D^t)=\begin{pmatrix}-2&0\\0&6\end{pmatrix}$.
La primera fila de $AX$ es la segunda de $X$ y la segunda
es el doble de la primera de $X$. Igualando filas:
\[
X=\boxed{\begin{pmatrix}0&3\\-2&0\end{pmatrix}}.
\]
Al sustituir, se obtiene el producto requerido.

\textbf{b)} Evaluamos la función lineal en los vértices:
\[
\begin{array}{c|rrrrr}
 &A&B&C&D&E\\\hline
F&2&16&24&23&2
\end{array}
\]
El máximo es $\boxed{24}$, sólo en $C=(4,0)$.
El mínimo es $\boxed{2}$. Por linealidad lo alcanza todo el
segmento $AE$. Cualquier punto que use con peso positivo
otro vértice en una combinación convexa tiene un valor mayor,
así que el segmento reúne todos los mínimos.
